\section{A uniform bound for the support of a pencil}
\label{sec:support}

Fix a normal integral projective model $X\subset\Pj_l^n$ of $L/l$.
For $h\in V_L$, define
\[
 \supp_X(h)=\{D\subset X:\ D\text{ is a prime divisor and }
                         \ord_D(h)\ne0\text{ in }\Lambda\}.
\]
Its degree is
\[
 s_X(h)=\sum_{D\in\supp_X(h)}\deg(D).
\]
The orders modulo $\ell$, and hence this finite set and its degree, depend
only on the class $h$.  We will prove that, for each $t\in K\setminus k$,
the integers $s_X(\Theta([t-a]))$ are bounded independently of $a\in k$.
The method is to realize any finite selection of these classes on a
residual curve of bounded genus.  The point correspondence from
Section~\ref{sec:residual-curves} then turns the desired bound into a
Riemann--Hurwitz estimate.  The bounded-genus and Hurwitz method appears
in \cite[Section~6, proof of Proposition~6.1]{BT2011}; we adapt the
bounded-support construction of \cite[Sections~4--6]{InternalBP} to
finite coefficients.

\subsection{Finite Kummer tests}

We first record why only finitely many classes can disappear in a
finitely generated extension.

\begin{lemma}\label{lem:kummer-kernel}
Let $E/F$ be a finitely generated field extension, where
$\charac F\ne\ell$ and $\mu_\ell\subset F$.  The kernel of
$V_F\longrightarrow V_E$ is finite dimensional over $\Lambda$.
\end{lemma}

\begin{proof}
Let $F'$ be the relative algebraic closure of $F$ in $E$.
Choose a transcendence basis ${\bf x}$ for $E/F$ such that
$[E:F({\bf x})]=N<\infty$.  For every finite extension $F_1/F$
contained in $F'$, the tuple ${\bf x}$ remains algebraically independent
over $F_1$, and
\[
 [F_1:F]=[F_1({\bf x}):F({\bf x})]\le N.
\]
The finite subextensions of $F'/F$ form a directed system with bounded
degrees.  One of them has maximal degree and contains all the others;
thus $F'/F$ is finite.

If $r$ independent classes in $V_F$ vanish in $V_E$, choose
representatives $f_1,\ldots,f_r\in F^\times$ and their $\ell$-th roots
in $E$.  Those roots belong to $F'$, and Kummer theory gives
\[
 [F(f_1^{1/\ell},\ldots,f_r^{1/\ell}):F]=\ell^r.
\]
Hence $\ell^r\le [F':F]$, which bounds the kernel dimension.
\end{proof}

The following geometric observation explains why a single smooth curve
can preserve any specified finite Kummer extension.  Its irreducibility
assertion is useful because the cover can have arbitrarily large degree.

\begin{lemma}\label{lem:support-bertini}
Let $U$ be a smooth integral locally closed subvariety of projective space
over an algebraically closed field, of dimension at least two.  Let
$Y\to U$ be a finite \'{e}tale morphism with $Y$ integral.
For general hyperplanes $H$, both $U\cap H$ and
$Y\times_U(U\cap H)$ are nonempty, smooth, and integral.
\end{lemma}

\begin{proof}
Smoothness of $U\cap H$ is Bertini's smoothness theorem for an embedded
smooth variety, valid in every characteristic
\cite[Lemma~33.47.3, Tag~0FD6]{Stacks}.  It implies smoothness
upstairs because the cover is \'{e}tale.  We give the irreducibility
argument, applied to either $Y\to\Pj^n$ or $U\to\Pj^n$.

Write $f:Y\to\Pj^n$ for the resulting quasi-finite morphism and
$a=\dim Y\ge2$.  Let $\mathcal H=(\Pj^n)^\vee$ be the space of
hyperplanes.  In the incidence of triples $(y_1,y_2,H)$ with
$f(y_1),f(y_2)\in H$, the locus $f(y_1)\ne f(y_2)$ is a projective
space bundle with fiber $\Pj^{n-2}$ over a nonempty open of $Y\times Y$.
It is therefore irreducible of dimension $2a+n-2$.
The locus of pairs with equal images has dimension at most $a$, by
quasi-finiteness; its hyperplane incidence has dimension at most
$a+n-1<2a+n-2$.

For a general hyperplane, $f^{-1}(H)$ is nonempty and has pure dimension
$a-1$.  Nonemptiness follows because the image of $f$ contains an open
of its $a$-dimensional closure; purity follows from the principal ideal
theorem on the integral variety $Y$.  Thus every component of the square
of the generic hyperplane section has dimension $2a-2$ over the function
field of $\mathcal H$.  No such component can be contained in the
equal-image incidence, whose total dimension is too small.  The square
of the generic section is consequently irreducible.

This implies geometric irreducibility of the generic section.  Indeed,
if it had more than one geometric component, ordered pairs of points
lying on the same component and on different components would give two
distinct Galois-invariant unions of components of its square.
Its already established generic smoothness excludes nonreducedness.
Geometric integrality spreads to a nonempty open of $\mathcal H$,
which proves the assertion for general hyperplanes.
\end{proof}

\begin{proposition}[Finite curve test]\label{prop:finite-curve-test}
There is an integer $G_X\ge0$, depending only on the embedded model $X$,
with the following property.  Given $g_1,\ldots,g_m\in L^\times$ whose
classes in $V_L$ are independent, there is a quasi-prime
$(d-1)$-divisor $w$ of $L/l$ such that:
\begin{enumerate}
\item $lw$ is algebraically closed and algebraic over a finite field of
characteristic different from $\ell$, and $Lw$ is the function field of a
smooth projective curve $C_w/lw$ of genus at most $G_X$;
\item the $g_j$ are $w$-units and their residue classes are independent
in $V_{Lw}$;
\item for every $j$,
\[
 \#\{x\in C_w(lw):\ord_x(\overline{g_j})\not\equiv0\pmod\ell\}
       =s_X([g_j]).
\]
\end{enumerate}
\end{proposition}

\begin{proof}
We first construct a curve over $l$, then specialize its finite defining
data, and finally realize that specialization by a valuation of $L$.

\emph{A curve preserving orders and independence.}
Let $D_1,\ldots,D_q$ be the prime divisors occurring in the divisors of
the $g_j$.  There is a closed subset $B\subset X$ of codimension at least
two such that $X\setminus B$ is smooth, the $D_i\setminus B$ are smooth
pairwise disjoint Cartier divisors, and near each $D_i\setminus B$ every
$g_j$ is a power of a local equation of $D_i$ times a unit.
Outside their union the functions are units.  To obtain $B$, remove the
singular locus of the normal variety $X$, the singular and non-Cartier
loci of the $D_i$, their pairwise intersections, and the exceptional
loci of these local expressions.  Each has codimension at least two.

On the smooth open
$U=X\setminus(B\cup D_1\cup\cdots\cup D_q)$, the simultaneous root
cover defined by the equations $z_j^\ell=g_j$ is finite \'{e}tale.
Kummer theory and independence say that it is integral of degree
$\ell^m$.  Successively choose $d-1$ general hyperplanes.  Apply
Lemma~\ref{lem:support-bertini} to the cover on each successive open
section, and Bertini smoothness to the base and its specified divisors.
We obtain a smooth integral projective curve $C\subset X$ that avoids
$B$, meets every $D_i$ transversely in exactly $\deg D_i$ distinct
points, and has an integral root cover over $C\cap U$.
The latter cover still has degree $\ell^m$, so the restricted classes
of the $g_j$ are independent.  At a point of $C\cap D_i$ their orders
are exactly their orders along $D_i$.  These are all their zeros and
poles on $C$, proving the required equality of support counts over $l$.

The genus of this curve is independent of the chosen functions.
At each stage we may also require the hyperplane to avoid the associated
points of the preceding scheme section.  Multiplication by its equation
is then injective, so the Hilbert polynomial of a section is the first
difference of the preceding Hilbert polynomial.  The final scheme
section avoids $B$, and on $X\setminus B$ the successive sections are
smooth by Bertini.  Hence the final scheme section is reduced and equals
$C$, even if intermediate sections had embedded components supported
in $B$.
Its Hilbert polynomial, and thus its genus, depend only on the Hilbert
polynomial of the embedded $X$.  Denote this genus by $G_X$.
This argument does not require $X$ to be Cohen--Macaulay.

\emph{Specialization of the finite data.}
Choose a finitely generated subring $R\subset l$ in which $\ell$ is
invertible and over which the preceding data are defined.  Enlarge $R$
by finitely many elements and localize it as follows.  The smooth open
of $X$ containing $C$, the flag sections in that open, and their
inclusions descend with their dimensions, smoothness, and geometric
integrality preserved in every geometric fiber.  The final curve
descends as a smooth projective curve of genus $G_X$.
Include the coordinates of its finitely many zeros and poles and the
local parameter--unit expressions for the functions.  After shrinking
$\Spec R$, they give disjoint point sections with the same orders and
no other zeros or poles.  The root cover over the complement of those
sections remains finite \'{e}tale of degree $\ell^m$, with geometrically
integral fibers.  Thus independence and the support counts persist in
every geometric fiber under consideration.

Here the spreading statements have their usual finite-presentation
meaning: finitely many schemes, morphisms, functions, and identities
descend after adjoining finitely many coefficients
\cite[Tags~01ZM, 0C0C, 081F]{Stacks}; smoothness and
flatness hold after shrinking; and geometric integrality at the generic
point persists on an open \cite[Tags~0578, 0559]{Stacks}.
The Hilbert polynomial in the projective
flat family of curves is constant; its constant term is the fiberwise
Euler characteristic, locally constant in this proper flat family
\cite[Lemma~36.32.2, Tag~0B9T]{Stacks}.  For the orders, the identities
$g_j=u\pi^e$ on finitely many neighborhoods, with $u$ invertible and
$\pi$ a parameter for the relevant section, preserve the integer $e$.
On the complement of these neighborhoods the functions and their
inverses are regular.  Properness of the curve ensures that these
finitely many open conditions still cover every fiber after shrinking.

We also include a presentation of the function field, so that the
specialized variety will be the residue field of a valuation.
Choose a separating transcendence basis $x_1,\ldots,x_d$ for $L/l$ and
a primitive element $y$ for the finite separable extension
$L/l(x_1,\ldots,x_d)$.  Write its monic minimal polynomial as
$Q(Y)\in l({\bf x})[Y]$.  After inverting a polynomial in ${\bf x}$,
this gives a finite integral model over an open of affine $d$-space.
Descend that model, a common dense open with the model already chosen,
and the rational expressions for all the $g_j$.
By geometric integrality and further localization of $R$, the reduced
polynomial in every geometric fiber is irreducible of the same degree
as $Q$ and gives that fiber's function field.  All denominators in the
coefficients and in the expressions being used remain nonzero, and the
$g_j$ agree with the specified nonzero functions there.

\emph{Realization by a valuation.}
Choose a closed point $s\in\Spec R$.  Its residue field is finite, of
characteristic different from $\ell$.  A valuation of $\operatorname{Frac}R$
dominating $R_s$ extends to $l$; equivalently, one may directly choose a
valuation ring of $l$ dominating $R_s$
\cite[Lemma~10.50.2, Tag~00IA]{Stacks}.
Write its residue field as $\lambda_0$.
Since $l$ is algebraically closed, so is $\lambda_0$.
We may arrange that the ultimate residue field is algebraic over the
finite field $\kappa(s)$.  Indeed let $b$ be the algebraic closure of
$\kappa(s)$ inside $\lambda_0$, and choose a transcendence basis
$\mathcal T$ for $\lambda_0/b$.  On $b(\mathcal T)$ give the basis
elements rationally independent values in an ordered free abelian group
with basis $\mathcal T$ and give $b^\times$ value zero.
Every polynomial has a unique term of least value, so this defines a
valuation with residue $b$.  Extend it to the algebraic extension
$\lambda_0/b(\mathcal T)$.  Its residue extension is algebraic and thus
still equals $b$.  Composing valuations gives a valuation $u$ of $l$
with residue $\lambda=b$.  Its center on $R$ remains $s$, since the
second valuation is trivial on the finite field $\kappa(s)$.
An ordered free abelian group exists for any cardinality of
$\mathcal T$, so this construction imposes no cardinality condition on $l$.

Give $l({\bf x})$ the Gauss extension of $u$: the values of the $x_i$
are zero and their residues are algebraically independent over
$\lambda$.  Prolong it to $L$, obtaining $w_X$.
The coefficients of $Q$ are integral at the Gauss valuation and its
reduction is the specified irreducible polynomial over
$\lambda(\overline{\bf x})$.  Since $Q$ is monic, $y$ is integral;
its residue therefore has degree $\deg Q$.  The fundamental inequality
for a finite extension of valued fields now forces
\[
 w_XL=ul,\qquad
 Lw_X=\lambda(\overline{\bf x},\overline y).
\]
Thus $w_X$ is a constant reduction with precisely the specialized
function field as residue, and each $g_j$ has the prescribed residue.

On $Lw_X$, take the successive true divisor valuations of the
specialized smooth flag, ending at its curve.  Compose them with $w_X$.
Write $w_1$ for the first composition.  It is quasi-divisorial: its value group $\Gamma$
has a convex kernel $\Z$ over $w_XL$, while the constant values map
isomorphically onto $w_XL=ul$.  Consequently
$\Gamma/w_1l\cong\Z$.  Every nonzero convex subgroup of $\Gamma$
contains this innermost $\Z$, so any proper coarsening has value group
generated by constant values and loses the relative rank-one
contribution.  This proves the required minimality.
The remaining steps are true divisor steps.  Their composition $w$ is
therefore a quasi-prime $(d-1)$-divisor with residue the specialized
curve field.  All $g_j$ are units at the successive generic points.
The genus, independence, and support counts already preserved in the
specialization give the three conclusions.
\end{proof}

\subsection{Five fibers control the degree}

We now transfer the finite tests across $\Theta$.
The curve used for a test is allowed to depend on the tested classes;
its genus bound depends only on $X$.

\begin{proposition}\label{prop:uniform-support}
For every $t\in K\setminus k$,
\[
 \sup_{a\in k}s_X\bigl(\Theta([t-a])\bigr)<\infty.
\]
\end{proposition}

\begin{proof}
Put $h_a=\Theta([t-a])$.  The classes $[t-a]$, $a\in k$, are
independent in $V_{k(t)}$, as is seen from their orders at the finite
points of the projective line.
Lemma~\ref{lem:kummer-kernel} shows that the kernel of
$V_{k(t)}\to V_K$ is finite dimensional.  A basis for its intersection
with the span of these classes involves only finitely many indices.
After removing those indices, we have an infinite set $A\subset k$
such that the classes $[t-a]$, $a\in A$, are independent in $V_K$.

Choose five distinct anchors $a_1,\ldots,a_5\in A$ and put
$N_i=s_X(h_{a_i})$.  If $\charac k=0$, choose the anchors to include
$b_0,b_0+1,b_0+\ell$; this is possible because $k\setminus A$ is finite.
Fix one further $a\in A$ distinct from the anchors and put $a_6=a$.
Choose representatives $g_1,\ldots,g_6\in L^\times$ of
$h_{a_1},\ldots,h_{a_6}$, and apply
Proposition~\ref{prop:finite-curve-test} to these representatives.
Let $w$ be the resulting valuation and $C_w$ its residual curve.
Their residue classes are independent, so their evaluations on $D_w$
are independent by Lemma~\ref{lem:residue-duality}.
Let $v$ be the quasi-prime $(d-1)$-divisor of $K$ corresponding to $w$
under Corollary~\ref{cor:flag-matching}.
The evaluations of the six functions $u_j=t-a_j$, $1\le j\le6$,
on $D_v$ are independent.

We first normalize these functions so that their residues are a pencil
on $Kv$.  We claim that $v(k(t))=vk$.
Otherwise the nonzero relative value group of $k(t)/k$ has rational
rank one, and the valuation transcendence-degree inequality forces
$k(t)v=kv$.  Equality in that inequality makes the relative value
group cyclic.  Any element with value in $vk$ can then be multiplied
by a constant to become a unit, and its residue can be removed by a
further constant.  Its class consequently vanishes on $D_v$.
Thus evaluation of $k(t)^\times$ on $D_v$ factors through the cyclic
relative value group and has dimension at most one, a contradiction.

The values of all six $u_j$ are equal.  If $v(u_i)>v(u_j)$, then
$u_j/(a_i-a_j)$ is a principal unit, contradicting its nonzero
evaluation on $D_v$.  Write $\gamma$ for the common value.
For distinct indices, $v(a_i-a_j)\ge\gamma$; a strict inequality
would make $u_i/u_j$ a principal unit and give equal evaluations.
Hence every difference $a_i-a_j$ has value exactly $\gamma$.
In positive characteristic the residue characteristic equals
$\charac k\ne\ell$.  In characteristic zero our three specified
anchors give $v(1)=v(\ell)=\gamma=0$, so again
$\charac(kv)\ne\ell$.

Choose $b\in k^\times$ with $v(b)=\gamma$ and set
\[
 \tau_0=\overline{(t-a_1)/b},\qquad
 \beta_j=\overline{(a_j-a_1)/b}.
\]
The $\beta_j\in kv$ are distinct and
$\overline{u_j/b}=\tau_0-\beta_j$.
These residue classes remain independent: the functions $u_j/b$ are
units, their evaluations factor through $D_v/I_v$, and multiplication
by $b$ changes no class because constants are $\ell$-divisible.
In particular $\tau_0\notin kv$.
More precisely, the dual of the induced isomorphism
$D_w/I_w\longrightarrow D_v/I_v$ carries
$[\tau_0-\beta_j]\in V_{Kv}$ to $[\overline{g_j}]\in V_{Lw}$.
Indeed their evaluations on corresponding residual characters are the
evaluations of $[u_j]$ and $\Theta([u_j])=[g_j]$ on the original
decomposition spaces.  This identifies the individual tested classes,
as well as preserving their independence.

Proposition~\ref{prop:point-recognition} says that $kv$, like $lw$,
is algebraic over a finite field, and that the complete point families
on the two residual curves correspond.  Let $C_v$ be the smooth
projective curve with function field $Kv$.
Both residue characteristics are different from $\ell$; by
Lemma~\ref{lem:curve-orders}, the dimensions of the common point-order
kernels are twice the genera.  The residual isomorphism preserves
these kernels, so $g(C_v)=g(C_w)\le G_X$.
It also preserves the number of nonzero point orders of each tested
class.  For the five anchor classes these numbers are $N_1,\ldots,N_5$.

Write $p=\charac(kv)>0$ and $\tau_0=\tau^{p^r}$ with $r$ maximal.
Such an $r$ exists because a nonconstant function has a nonzero point
order, and $p^r$ must divide that fixed nonzero integer.
The function $\tau$ defines a separable morphism
$C_v\to\Pj_{kv}^1$; write $n'$ for its degree.
Multiplication of orders by $p^r$ does not change their vanishing
modulo $\ell$.

Consider the fiber of this morphism over $\beta_i^{1/p^r}$.
At most $N_i$ points in the fiber have ramification index not divisible
by $\ell$.  The other points have index at least $\ell$, so there are
at most $n'/\ell$ of them.  If $r_i$ is the total number of points
in the fiber, then
$r_i\le N_i+n'/\ell$.
Since the sum of the ramification indices in a fiber is $n'$, and
the different exponent at a point is at least its index minus one,
the fiber contributes at least
$n'-r_i\ge(1-1/\ell)n'-N_i$ to the different.
The five fibers are disjoint.  Riemann--Hurwitz, with its different term
in arbitrary characteristic \cite[Lemma~53.12.4, Tag~0C1F]{Stacks},
therefore yields
\[
 2g(C_v)-2+2n'
   \ge 5(1-1/\ell)n'-\sum_{i=1}^5N_i,
\]
or
\begin{equation}\label{eq:support-five-fibers}
 \bigl(5(1-1/\ell)-2\bigr)n'
       \le 2G_X-2+\sum_{i=1}^5N_i.
\end{equation}
The coefficient on the left is positive for every prime $\ell$,
including $\ell=2$.  Hence $n'$ has a bound depending only on the five
anchors and $X$.

For the additional parameter, the zeros and poles of
$\tau_0-\beta_6$ lie in two fibers of this same morphism.
It has at most $2n'$ nonzero point orders modulo $\ell$.
The point correspondence and the finite curve test identify this
number with $s_X(h_a)$.  This proves a uniform bound for all
non-anchor members of $A$.  Including the anchors and the finitely
many parameters outside $A$ completes the proof.
\end{proof}
